% !TeX root = ../main.tex
\subsection*{2. Схемы и расчет коэффициентов отражения аттенюаторов и делителя мощности}

\textbf{Условие:} \\
Составить схемы резистивного П-образного и Т-образного аттенюаторов, а также резистивного делителя мощности 1:2. Выбрать коэффициент ослабления аттенюатора ($A = 10\text{ дБ}$). Рассчитать для каждой схемы коэффициенты отражения двухполюсника «аттенюатор--нагрузка» при волновом сопротивлении тракта $Z_0 = 50\text{ Ом}$ для двух вариантов нагрузки:
\begin{itemize}
    \item $Z_L = 50\text{ Ом}$ (для делителя мощности второй порт нагружен на $50\text{ Ом}$);
    \item $Z_L = 25\text{ Ом}$ (для делителя мощности второй порт нагружен на $50\text{ Ом}$).
\end{itemize}

\vspace{0.3cm}
\textbf{1. Схемы и расчет номиналов резисторов ($A = 10\text{ дБ}$, $Z_0 = 50\text{ Ом}$):}

Коэффициент ослабления по напряжению:
$$K = 10^{\frac{A}{20}} = 10^{\frac{10}{20}} \approx 3{,}162$$

\begin{itemize}
    \item \textbf{Т-образный аттенюатор:} состоит из двух последовательных резисторов $R_1$ и одного параллельного $R_2$:
    $$R_1 = Z_0 \frac{K - 1}{K + 1} = 50 \cdot \frac{3{,}162 - 1}{3{,}162 + 1} \approx 25{,}98\text{ Ом}$$
    $$R_2 = Z_0 \frac{2K}{K^2 - 1} = 50 \cdot \frac{2 \cdot 3{,}162}{3{,}162^2 - 1} \approx 35{,}14\text{ Ом}$$

    \item \textbf{П-образный аттенюатор:} состоит из одного последовательного резистора $R_1$ и двух параллельных $R_2$:
    $$R_1 = Z_0 \frac{K^2 - 1}{2K} = 50 \cdot \frac{3{,}162^2 - 1}{2 \cdot 3{,}162} \approx 71{,}15\text{ Ом}$$
    $$R_2 = Z_0 \frac{K + 1}{K - 1} = 50 \cdot \frac{3{,}162 + 1}{3{,}162 - 1} \approx 96{,}25\text{ Ом}$$

    \item \textbf{Резистивный делитель мощности 1:2:} симметричная схема на двух резисторах $R_1 = R_2 = 50\text{ Ом}$, включенных последовательно в каждое выходное плечо.
\end{itemize}

\vspace{0.3cm}
\textbf{2. Расчет коэффициентов отражения для согласованного аттенюатора ($A = 10\text{ дБ}$):}

Для любого согласованного симметричного аттенюатора ($S_{11} = S_{22} = 0$, $S_{21} = S_{12} = 1/K$) входной коэффициент отражения выражается через коэффициент отражения нагрузки $\Gamma_L$:
$$\Gamma_{\text{вх}} = S_{11} + \frac{S_{21} S_{12} \Gamma_L}{1 - S_{22} \Gamma_L} = \frac{1}{K^2} \Gamma_L = 10^{-\frac{2A}{20}} \Gamma_L = 0{,}1 \cdot \Gamma_L$$
В логарифмическом масштабе: $\Gamma_{\text{вх [дБ]}} = \Gamma_{L\text{ [дБ]}} - 2A$.

\begin{itemize}
    \item \textbf{При нагрузке $Z_L = 50\text{ Ом}$:}
    $$\Gamma_L = \frac{Z_L - Z_0}{Z_L + Z_0} = \frac{50 - 50}{50 + 50} = 0 \quad (-\infty\text{ дБ})$$
    $$\Gamma_{\text{вх}} = 0 \quad (-\infty\text{ дБ}), \quad \text{КСВН} = 1$$

    \item \textbf{При нагрузке $Z_L = 25\text{ Ом}$:}
    $$\Gamma_L = \frac{25 - 50}{25 + 50} = -\frac{1}{3} \approx -0{,}333 \quad (\Gamma_{L\text{ [дБ]}} = 20 \lg(1/3) \approx -9{,}54\text{ дБ})$$
    $$\Gamma_{\text{вх}} = 0{,}1 \cdot \left(-\frac{1}{3}\right) \approx -0{,}0333$$
    $$\Gamma_{\text{вх [дБ]}} = -9{,}54 - 2 \cdot 10 = -29{,}54\text{ дБ}$$
    $$\text{КСВН} = \frac{1 + |\Gamma_{\text{вх}}|}{1 - |\Gamma_{\text{вх}}|} = \frac{1 + 0{,}0333}{1 - 0{,}0333} \approx 1{,}069$$
\end{itemize}

\vspace{0.3cm}
\textbf{3. Расчет коэффициентов отражения для резистивного делителя мощности 1:2:}

Входное сопротивление делителя определяется параллельным соединением двух ветвей: $Z_{\text{вх}} = (R_1 + Z_{L1}) \parallel (R_2 + Z_{L2})$, где второе плечо всегда нагружено на $Z_{L2} = 50\text{ Ом}$, следовательно, ветвь 2 имеет сопротивление $R_2 + Z_{L2} = 50 + 50 = 100\text{ Ом}$.

\begin{itemize}
    \item \textbf{При нагрузке первого плеча $Z_{L1} = 50\text{ Ом}$:}
    $$R_1 + Z_{L1} = 50 + 50 = 100\text{ Ом}$$
    $$Z_{\text{вх}} = \frac{100 \cdot 100}{100 + 100} = 50\text{ Ом}$$
    $$\Gamma_{\text{вх}} = \frac{50 - 50}{50 + 50} = 0 \quad (-\infty\text{ дБ}), \quad \text{КСВН} = 1$$

    \item \textbf{При нагрузке первого плеча $Z_{L1} = 25\text{ Ом}$:}
    $$R_1 + Z_{L1} = 50 + 25 = 75\text{ Ом}$$
    $$Z_{\text{вх}} = \frac{75 \cdot 100}{75 + 100} = \frac{7500}{175} \approx 42{,}86\text{ Ом}$$
    $$\Gamma_{\text{вх}} = \frac{42{,}86 - 50}{42{,}86 + 50} \approx -0{,}0769$$
    $$\Gamma_{\text{вх [дБ]}} = 20 \lg(0{,}0769) \approx -22{,}28\text{ дБ}$$
    $$\text{КСВН} = \frac{1 + 0{,}0769}{1 - 0{,}0769} \approx 1{,}167$$
\end{itemize}

\vspace{0.3cm}
\textbf{Сводная таблица результатов:}
\begin{center}
\begin{tabular}{|l|c|c|c|c|}
\hline
Схема & Нагрузка $Z_L$, Ом & $\Gamma_{\text{вх}}$ (лин.) & $\Gamma_{\text{вх}}$, дБ & КСВН \\
\hline
Аттенюатор 10 дБ (П/Т) & 50 & 0 & $-\infty$ & 1{,}000 \\
Аттенюатор 10 дБ (П/Т) & 25 & $-0{,}0333$ & $-29{,}54$ & 1{,}069 \\
\hline
Делитель мощности 1:2 & 50 & 0 & $-\infty$ & 1{,}000 \\
Делитель мощности 1:2 & 25 & $-0{,}0769$ & $-22{,}28$ & 1{,}167 \\
\hline
\end{tabular}
\end{center}

\vspace{0.3cm}
\textbf{Влияние выбора другого коэффициента ослабления:}
\begin{itemize}
    \item Отраженная от нагрузки волна проходит через аттенюатор дважды (в прямом и обратном направлении), поэтому входной коэффициент отражения всегда улучшается ровно на удвоенное значение затухания: $\Gamma_{\text{вх [дБ]}} = \Gamma_{L\text{ [дБ]}} - 2A$.
    \item При выборе аттенюатора $3\text{ дБ}$: $\Gamma_{\text{вх [дБ]}} = -9{,}54 - 6 = -15{,}54\text{ дБ}$ ($\text{КСВН} \approx 1{,}41$).
    \item При выборе аттенюатора $6\text{ дБ}$: $\Gamma_{\text{вх [дБ]}} = -9{,}54 - 12 = -21{,}54\text{ дБ}$ ($\text{КСВН} \approx 1{,[118;1:3u}18$).
    \item Чем больше ослабление аттенюатора, тем ближе входное сопротивление к $50\text{ Ом}$ независимо от величины несогласованной нагрузки.
\end{itemize}
